% =========================
% 7_discussion.tex --- structural limits of the design, and what is not yet claimed
% =========================
\section{Discussion and Limitations}
\label{sec:discussion}

\textbf{Same-configuration cells cannot test admissibility.}
The gate scores the query sensor and the evidence sensor against each candidate and combines them.
When enrollment and query come from the same configuration, both factors read the same description,
and the product reduces to a function of the candidate alone. A query-independent per-candidate
constant multiplying an evidence-derived vote can shift an argmax, but only toward whichever concept
the gate considers a priori more observable, regardless of what the retrieved evidence says. This is
a property of the rule, not of any particular fit, and it is why
Section~\ref{sec:eval_protocol} treats cross-configuration cells as a precondition rather than an
extension. It also means an evaluation restricted to same-configuration cells can measure the
mechanism's cost while being structurally unable to measure its benefit.

\textbf{The veto threshold is not tied to the distribution it gates.}
Admissibility below a fixed constant removes a row from a candidate's pool. That constant does not
move with the fitted gate, whose output scale depends on the encoder that produced the resolvability
table. Two failure modes follow. If the neutral value that unfamiliar text falls back to sits above
the threshold, abstention rescues every unfamiliar configuration and the veto never fires, leaving
only reweighting. If it sits below, an unfamiliar configuration is vetoed on every candidate and
contributes nothing. Nothing in the current design asserts an ordering between the two. Expressing
the threshold as a quantile of the fitted values, and asserting that ordering at fit time, would
remove the failure mode; we have not done so.

\textbf{Provenance is confounded with hardware in public corpora.}
Each public dataset generally uses one device model and one preprocessing chain, so a descriptor of
acquisition behavior is also a descriptor of which study a recording came from. We respond by
keeping the acquisition-statistics similarity term disabled in the reported retriever and by
measuring same-dataset retrieval directly rather than assuming it is benign. The confound is a
property of the available data and no estimator resolves it. Settling it needs a corpus with the
same hardware across studies, or the same study across hardware.

\textbf{Scope of the recognition problem.}
HALO classifies presegmented single-activity windows. It has no unknown output and does not perform
continuous segmentation or calibrated rejection, so a deployment that encounters an activity outside
its candidate set will still name one. Selective classification would require truth-absent episodes
and risk--coverage evaluation~\cite{selectiveclass}; we treat it as future work rather than claiming
it, because the evaluation artifacts that would support such a claim do not exist.

\textbf{No accuracy and no on-device claim.}
We report the design, the implementation, and the protocol, and no recognition accuracy. The
encoder is complete and its memory is built, but the gate has not been fitted against that memory,
so any number we printed would describe a different system than the one specified here. We likewise
make no claim about phone-side latency, energy, or export, because we have not built that path. The
design does bound what such a claim would look like: the encoder is 7.70\,M parameters and
enrollment appends rows rather than running a gradient step, so the runtime question is retrieval
over a bank rather than on-device training. The bank is the obvious pressure point --- 6.9\,GB is a
server-side artifact, and a phone-side deployment would need either a per-user subset or a
compressed index.

\textbf{Sealing the test roster.}
Seven evaluation datasets remain unconsumed, and Section~\ref{sec:pending_measurements} states the
criteria for unsealing them. This is a deliberate cost: it means this paper reports no external
performance number. We think it is the correct trade while the development criteria are unmet,
because a test number produced under a protocol that cannot express the mechanism would not be
evidence for the claim the paper makes.
